% wave12_b3_6d_hCS_E3_gen_rel.tex
% Frontier working note --- Wave 12 / B3
% Channel: Costello-Francis-Gwilliam 2026 + Cattaneo-Felder-Mnev precision.
% Mission: construct, from first principles, the 6d holomorphic Chern-Simons
% avatar of the E_3 algebra of observables, in generators and relations.
% Output scope: (i) classical BV action; (ii) BRST/BV differential;
% (iii) pre-factorization pre-cosheaf Obs^{cl}; (iv) quantum RG + Feynman
% coefficients producing Obs^{q}; (v) E_3 algebra-of-observables at a point.
% Five ATTACK <-> HEAL cycles.

\documentclass[11pt]{amsart}
\usepackage[T1]{fontenc}
\usepackage[utf8]{inputenc}
\usepackage{amsmath,amssymb,amsthm}
\usepackage{mathtools}
\usepackage[margin=1.1in]{geometry}
\usepackage{hyperref}

\providecommand{\C}{\mathbb{C}}
\providecommand{\R}{\mathbb{R}}
\providecommand{\Z}{\mathbb{Z}}
\providecommand{\N}{\mathbb{N}}
\providecommand{\Q}{\mathbb{Q}}
\providecommand{\dbar}{\bar{\partial}}
\providecommand{\g}{\mathfrak{g}}
\providecommand{\h}{\mathfrak{h}}
\providecommand{\Obs}{\mathrm{Obs}}
\providecommand{\Sym}{\mathrm{Sym}}
\providecommand{\End}{\mathrm{End}}
\providecommand{\Hom}{\mathrm{Hom}}
\providecommand{\id}{\mathrm{id}}
\providecommand{\gh}{\mathrm{gh}}
\providecommand{\af}{\mathrm{af}}
\providecommand{\cl}{\mathrm{cl}}
\providecommand{\qu}{\mathrm{q}}
\providecommand{\Fact}{\mathrm{Fact}}
\providecommand{\PreFact}{\mathrm{PreFact}}
\providecommand{\E}{E}
\providecommand{\Disk}{\mathrm{Disk}}
\providecommand{\cE}{\mathcal{E}}
\providecommand{\cA}{\mathcal{A}}
\providecommand{\cO}{\mathcal{O}}
\providecommand{\cU}{\mathcal{U}}
\providecommand{\cF}{\mathcal{F}}
\providecommand{\cP}{\mathcal{P}}
\providecommand{\cS}{\mathcal{S}}
\providecommand{\kappach}{\kappa_{\mathrm{ch}}}
\providecommand{\kappaBKM}{\kappa_{\mathrm{BKM}}}
\providecommand{\hCS}{\mathrm{hCS}}

\theoremstyle{plain}
\newtheorem{theorem}{Theorem}[section]
\newtheorem{proposition}[theorem]{Proposition}
\newtheorem{lemma}[theorem]{Lemma}
\newtheorem{corollary}[theorem]{Corollary}
\theoremstyle{definition}
\newtheorem{definition}[theorem]{Definition}
\newtheorem{construction}[theorem]{Construction}
\theoremstyle{remark}
\newtheorem{remark}[theorem]{Remark}
\newtheorem*{attack}{ATTACK}
\newtheorem*{heal}{HEAL}

\title{Six real dimensions of holomorphic Chern--Simons on $\C^3$:\\
BV action, pre-factorization observables, and the $\E_3$ avatar at a point}
\author{Wave 12 / B3 --- Costello--Francis--Gwilliam channel}
\date{2026-04-22}

\begin{document}
\maketitle

\begin{abstract}
Holomorphic Chern--Simons on a Calabi--Yau threefold $X$ with
$\dim_\C X = 3$, i.e.\ $\dim_\R X = 6$, has classical action
$S_\cl(\cA) = \tfrac12\int_X \Omega \wedge \langle \cA, \dbar \cA\rangle +
\tfrac16\int_X \Omega \wedge \langle \cA, [\cA, \cA]\rangle$ on the total
Dolbeault complex $\cA \in \Omega^{0,\ast}(X, \g)$. We construct:
(i) the BV complex with $(-1)$-shifted symplectic pairing; (ii) the BRST
differential $Q = \dbar + [-, -]$; (iii) the pre-factorization pre-cosheaf
$\Obs^\cl$ of classical observables; (iv) the Costello renormalised
quantum observables $\Obs^\qu$ with Bochner--Martinelli propagator and
one-loop Feynman-integral coefficients; (v) the $\E_3$-algebra structure
on $\Obs^\qu_{(0)}$ at a point, equivalent by Costello--Francis--Gwilliam
2026 to the universal enveloping $\E_3$-algebra of $\g[2]$, the
degree-$2$ shift of $\g$. Scope discipline: ``$6$d holomorphic
Chern--Simons'' names the $6$-real-dimensional ambient, not Costello's
$4$d+$2$d twist; the two are distinct theories.
\end{abstract}

\tableofcontents

%% =====================================================================
\section{Scope: what $6$d holomorphic Chern--Simons names here}
%% =====================================================================

\begin{attack}[Cycle 1, scope]
``Holomorphic Chern--Simons is 3-dimensional: $\dbar$-Chern--Simons on a
Calabi--Yau threefold is a functional $S(\cA) = \int \Omega \wedge
(\cA\,\dbar\cA + \tfrac13\cA[\cA,\cA])$, with $\cA$ a $(0,1)$-form on a
3-complex-dimensional base. Costello's 4d Chern--Simons adds a 2d chiral
sector; their sum is 6d holomorphic. Which 6d Chern--Simons is the
``$\E_3$ avatar at a point''?''
\end{attack}

\begin{heal}[Cycle 1 resolution]
Two distinct theories share the phrase ``holomorphic Chern--Simons'':
\begin{enumerate}
\item[(A)] Witten's hCS on a Calabi--Yau 3-fold $X$: the gauge
field is a $(0,1)$-form $A$, the action is Witten 1992
\cite{Witten1992CS}:
\[
S_\hCS(A) \;=\; \int_X \Omega \wedge \Bigl(\tfrac12 A\,\dbar A +
\tfrac13 A\,[A, A]\Bigr),
\qquad A \in \Omega^{0,1}(X, \g).
\]
Here $\dim_\R X = 6$; the theory is ``$6$d'' in real dimension,
``$3$d holomorphic'' in complex dimension.
\item[(B)] Costello's $4$d holomorphic Chern--Simons: the gauge
field is a $(0,1)$-form on $\Sigma \times C$ with $\Sigma$ complex
surface and $C$ complex curve; the action integrates against
$\omega_C \wedge \omega_\Sigma$ \cite{Costello2013}.
\item[(C)] Costello--Francis--Gwilliam ``$6$d'' or ``higher-dimensional''
holomorphic Chern--Simons: the 6-real-dimensional ambient with
$(0, \ast)$-valued field, BV-promoted per this note.
\end{enumerate}
The present construction is (A) with the BV extension that
promotes $A \in \Omega^{0,1}$ to the total $\cA \in \Omega^{0, \ast}$
via Costello--Gwilliam Vol.\ II \S5.\ The $\E_3$ avatar arises
from this theory, not from (B). In a slight abuse of convention
we call this theory ``$6$d hCS on $X$''; ``$6$d'' specifies the
real dimension of the ambient, not the form degree of the gauge
field.
\end{heal}

\begin{remark}[Distinction from Costello's 4d hCS]\label{rem:not-4d}
Costello's 4d hCS on $\Sigma \times C$ produces, on boundaries $\Sigma
\times \partial C$, the affine Yangian by Costello 2013; that theory has
bulk $\dim_\R = 4$. It is not subsumed by the 6d hCS here. The
$\E_3$ avatar in the present note arises only from the $6$-real-dimensional
case on $X$ with $\dim_\C X = 3$.
\end{remark}

\begin{construction}[Data and conventions]\label{cons:data}
Let $(X, \Omega)$ be a Calabi--Yau 3-fold, i.e.\ $\dim_\C X = 3$,
$K_X \simeq \cO_X$ trivialised by a nowhere-vanishing holomorphic
volume form $\Omega \in H^0(X, K_X)$. Let $\g$ be a Lie algebra
(not necessarily compact; for abelian passages we take $\g = \C$).
Fix an invariant non-degenerate pairing $\langle \cdot, \cdot\rangle
\colon \g \otimes \g \to \C$, either the Killing form when $\g$
semisimple, or the canonical pairing $\langle x, y\rangle = xy$ for
$\g = \C$. Write $\cE := \Omega^{0, \ast}(X, \g) = \bigoplus_{q=0}^{3}
\Omega^{0, q}(X, \g)$ with Dolbeault differential $\dbar$ of
$(0, q)$-to-$(0, q+1)$ degree.
\end{construction}

%% =====================================================================
\section{The classical BV action on $\Omega^{0, \ast}(X, \g)$}
%% =====================================================================

\begin{attack}[Cycle 2, BV extension]
``The hCS field $A$ is a $(0,1)$-form. Promoting to a total
$(0, \ast)$-form is a choice --- where does the BV master equation
$\{S, S\} = 0$ live? Is the total degree of the BV superfield $+1$, and
what is the $(-1)$-shifted symplectic pairing?''
\end{attack}

\begin{heal}[Cycle 2 resolution]
BV extension of a classical field theory with gauge symmetry
lives on $\cE = \Omega^{0, \ast}(X, \g)$ with a $(-1)$-shifted
cohomological degree, via the AKSZ construction
\cite{AKSZ1997}, \cite{CFL2000}. Write
\[
\cA \;=\; c + A_{0,1} + A_{0,2}^\ast + c_{0,3}^\ast, \qquad
\cA \in \Omega^{0, \ast}(X, \g),
\]
with
\[
\begin{array}{c|cccc}
 & c & A_{0,1} & A_{0,2}^\ast & c_{0,3}^\ast \\ \hline
(0, q) & 0 & 1 & 2 & 3 \\
\gh & +1 & 0 & -1 & -2 \\
\af & 0 & 0 & 1 & 2
\end{array}
\]
The total cohomological (BV) degree is
$\deg_{BV}(\cA) = (0, q) - 1$, placing $\cA$ homogeneously in
BV-degree $+1$ after a $1$-shift. The $(-1)$-shifted symplectic
pairing is $\omega_{BV}(\cA_1, \cA_2) = \int_X \Omega \wedge
\langle\cA_1, \cA_2\rangle$, pairing degrees $(0, q)$ and
$(0, 3-q)$ by Serre duality, and Koszul-dual $c \leftrightarrow
c^\ast_{0,3}$, $A_{0,1} \leftrightarrow A^\ast_{0,2}$.
\end{heal}

\begin{definition}[Classical BV action]\label{def:BV-action}
The classical BV action on $\cE = \Omega^{0, \ast}(X, \g)$ is
\[
S_\cl(\cA) \;=\; \int_X \Omega \wedge \Bigl(\tfrac12 \langle \cA, \dbar\cA\rangle
+ \tfrac16 \langle\cA, [\cA, \cA]\rangle\Bigr),
\qquad \cA \in \Omega^{0, \ast}(X, \g).
\]
This is the unique local functional on $\cE$ of BV-degree $0$ that
restricts to Witten's hCS on $\cA|_{(0,1)} = A$, and
whose Euler--Lagrange equation
\[
\dbar \cA + \tfrac12 [\cA, \cA] \;=\; 0
\]
is the Maurer--Cartan equation of the dgla
$(\Omega^{0, \ast}(X, \g), \dbar, [-, -])$.
\end{definition}

\begin{proposition}[Classical master equation]\label{prop:CME}
$\{S_\cl, S_\cl\} = 0$ is equivalent to
\[
\dbar^2 = 0, \qquad \text{Jacobi identity on } \g,
\qquad \dbar [\cA, \cA'] = [\dbar \cA, \cA'] + (-)^{|\cA|}[\cA, \dbar\cA'],
\]
i.e.\ to the dgla structure on $\cE$. All three are automatic. Hence
$\{S_\cl, S_\cl\} = 0$.
\end{proposition}

\begin{proof}
Compute $\{S, S\}_{BV} = \int \Omega \wedge \langle Q \cA, Q\cA\rangle$ with
$Q\cA = \dbar \cA + \tfrac12 [\cA, \cA]$. The cubic $\int \Omega \wedge
\langle[\cA, \cA], [\cA, \cA]\rangle$ vanishes by cyclic invariance of
$\langle\cdot, \cdot\rangle$ + Jacobi; the cross term vanishes by Leibniz
for $\dbar$ on $[-, -]$. See Costello--Gwilliam 2017 Vol.\ II Prop.\ 4.6.1
\cite{CGVol2}, applied to the dgla $\Omega^{0, \ast}(X, \g)$.
\end{proof}

\begin{definition}[BRST differential $Q$]\label{def:Q}
The BRST differential is
\[
Q \;=\; \dbar + \tfrac12[\cA, -].
\]
On linear functionals it acts as
$Q c = \tfrac12[c, c]$, $Q A_{0,1} = \dbar c + [c, A_{0,1}]$,
$Q A^\ast_{0,2} = \dbar A_{0,1} + \tfrac12[A_{0,1}, A_{0,1}] + [c, A^\ast_{0,2}]$,
$Q c^\ast_{0,3} = \dbar A^\ast_{0,2} + [A_{0,1}, A^\ast_{0,2}] +
[c, c^\ast_{0,3}]$. $Q^2 = 0$ is equivalent to the Jacobi identity on $\g$
and $\dbar^2 = 0$.
\end{definition}

\begin{remark}[Abelian linearisation]
At $\g = \C$, all brackets vanish and $Q = \dbar$. The BV complex becomes
$(\Omega^{0, \ast}(X), \dbar)$ with a $1$-shift in BV-degree. Its
cohomology is $H^{0, \ast}_{\dbar}(X) = H^\ast(X, \cO_X)$ with the
Hodge--Serre structure; for $X$ a compact Calabi--Yau 3-fold this has
dimensions $(h^{0,0}, h^{0,1}, h^{0,2}, h^{0,3}) = (1, 0, 0, 1)$.
\end{remark}

\begin{remark}[BRST cohomology = Lie cohomology]\label{rem:CE}
For general $\g$ and $X = \C^3$, the BRST cohomology of
$(\Omega^{0, \ast}(\C^3, \g), Q)$ is computable via the
Koszul--Chevalley--Eilenberg complex. Passing to derived
formal neighbourhoods at $0 \in \C^3$:
\[
H^\ast_{Q}\bigl(\Omega^{0, \ast}_0(\C^3, \g)\bigr) \;\simeq\;
H^\ast_{\mathrm{CE}}\bigl(\g[[z_1, z_2, z_3]], \cO_{\C^3, 0}\bigr),
\]
i.e.\ Lie-algebra cohomology of the algebra of $\g$-valued
power series, with coefficients in holomorphic functions on
$\C^3$. The $\E_3$-avatar identified in \S6 is an $\E_3$-refinement
of the coefficients in this CE complex.
\end{remark}

%% =====================================================================
\section{The pre-factorization pre-cosheaf of classical observables}
%% =====================================================================

\begin{attack}[Cycle 3, pre-factorization structure]
``What is the pre-factorization (pre-cosheaf) structure on classical
observables? Is $\Obs^\cl(U) = \Sym(\cE^\vee(U)[-1])$ actually a
pre-cosheaf, and does it satisfy the Weiss-cover descent condition?''
\end{attack}

\begin{heal}[Cycle 3 resolution]
Classical observables form a pre-cosheaf of cochain complexes on
open subsets of $X$, not a cosheaf in the sheaf-theoretic sense.
The Weiss (polydisk) cover property is the correct descent condition.
Costello--Gwilliam Vol.\ I \S3.1 \cite{CGVol1}:
\end{heal}

\begin{definition}[Classical observables pre-cosheaf]\label{def:Obs-cl}
For $U \subset X$ open, set
\[
\Obs^\cl(U) \;=\; \Sym_\C\bigl(\cE_c(U)^\vee[-1]\bigr)_{Q}
\]
with
\begin{itemize}
\item $\cE_c(U) = \Omega^{0, \ast}_c(U, \g) \subset \Omega^{0, \ast}(X, \g)$ the
compactly supported fields;
\item $\cE_c(U)^\vee$ its topological dual (distributional linear
functionals);
\item $[-1]$ a cohomological shift placing linear functionals
in BV-degree $-1$ so that $\Sym$ lands in non-negative BV-degrees;
\item $Q$ the differential induced by the BV differential on $\cE$.
\end{itemize}
For an inclusion $U \subset V$, extension by zero on $\cE_c(U) \subset
\cE_c(V)$ is covariant, giving a pre-cosheaf structure map
$\Obs^\cl(U) \to \Obs^\cl(V)$.
\end{definition}

\begin{proposition}[Pre-factorization structure]\label{prop:prefact}
The assignment $U \mapsto \Obs^\cl(U)$ is a pre-factorization algebra in
the sense of Costello--Gwilliam Vol.\ I Def.\ 3.1.4: for every
finite disjoint union $U_1 \sqcup \cdots \sqcup U_k \subset V$ of open
subsets of $X$, the multiplication
\[
m_{U_1, \ldots, U_k; V}\colon
\Obs^\cl(U_1) \otimes \cdots \otimes \Obs^\cl(U_k)
\;\longrightarrow\;
\Obs^\cl(V)
\]
sends $(\phi_1, \ldots, \phi_k)$ to the symmetric product
$\phi_1 \cdot \phi_2 \cdots \phi_k$ of extended-by-zero functionals.
Associativity follows from symmetry of $\Sym$; compatibility with $Q$
follows from $Q$ being a derivation of $\Sym$.
\end{proposition}

\begin{proposition}[Weiss descent, classical]\label{prop:Weiss-cl}
The pre-factorization pre-cosheaf $\Obs^\cl$ is a factorization algebra
in the sense of Costello--Gwilliam Vol.\ I Def.\ 6.1.1: Weiss-cover
descent holds, i.e.\ the homotopy colimit of $\Obs^\cl$ over any
Weiss cover of $U$ recovers $\Obs^\cl(U)$.
\end{proposition}

\begin{proof}
The Weiss condition for the functor $U \mapsto \Sym(\cE_c(U)^\vee[-1])$
reduces, by polynomial filtration, to the Weiss condition for the
underlying elliptic complex $\cE$. For $\cE = \Omega^{0, \ast}(-, \g)$
the Dolbeault complex, Weiss descent is classical fact,
Costello--Gwilliam Vol.\ I Thm.\ 6.5.3. The functor $\Sym(-[-1])$
preserves Weiss descent because it is a polynomial functor; $Q$ is
compatible because $\dbar$ is.
\end{proof}

\begin{construction}[$P_0$ bracket on $\Obs^\cl$]\label{cons:P0}
The $(-1)$-shifted symplectic pairing $\omega_{BV}$ induces, after
inverting off the diagonal, a $P_0$-bracket $\{-, -\}_\cl$ on
$\Obs^\cl(U)$ of cohomological degree $+1$:
\[
\{\phi, \psi\}_\cl \;=\; \int_{U \times U}\!\! \Omega_z \wedge \Omega_w \wedge
K_{BV}(z, w) \cdot \phi(z)\, \psi(w),
\]
with $K_{BV}$ the Bochner--Martinelli kernel introduced in \S4.
The bracket is graded-symmetric (odd Poisson) and satisfies graded
Leibniz; the Jacobi identity is the shadow of
$\{S_\cl, S_\cl\} = 0$ upon expansion. This is the
Costello--Gwilliam $P_0$-algebra structure of Vol.\ I \S5.4.
\end{construction}

%% =====================================================================
\section{Propagator and Feynman integrals}
%% =====================================================================

\begin{attack}[Cycle 4, renormalisation]
``Quantum observables are obtained from classical observables by a
Costello renormalisation group flow. This requires a propagator and
Feynman-diagram weights. On $\C^3$ the Bochner--Martinelli kernel is
the candidate; but its behaviour on the diagonal is singular. What are
the explicit 1-loop coefficients? Do they multiply $\hbar$ or $\hbar^2$?
Do they contribute counterterms that break the BV master equation?''
\end{attack}

\begin{heal}[Cycle 4 resolution]
The Bochner--Martinelli kernel provides the heat-kernel-regularised
propagator. Feynman coefficients at $n$-loop order multiply
$\hbar^n$. On $\C^3$ with abelian $\g = \C$ all loops above $1$ vanish
(Gaussian free theory); for non-abelian $\g$ the one-loop anomaly is
controlled by the second Chern class and can obstruct quantisation
unless $c_2(\g) = 0$ (Costello--Li 2015 \cite{CostelloLi}).
\end{heal}

\begin{proposition}[Bochner--Martinelli propagator]\label{prop:BM}
The inverse of $\dbar$ on $\Omega^{0, \ast}(\C^3)$ modulo harmonic
representatives is the Bochner--Martinelli kernel
\[
P(z, w) \;=\; \frac{2!}{(2\pi i)^3}
\sum_{k=1}^{3} \frac{(-1)^{k-1}\overline{(z_k - w_k)}\,
d\bar z_1 \wedge \widehat{d \bar z_k} \wedge d\bar z_2 \wedge \cdots}{|z - w|^6},
\]
a $(3, 2)$-current on $\C^3 \times \C^3 \setminus \Delta$ satisfying
$\dbar_w P(z, w) = \delta_\Delta$. This is the Schwartz-kernel of the
propagator pairing $A_{0,1}$ with $A^\ast_{0,2}$.
\end{proposition}

\begin{proof}
Bochner--Martinelli integral formula, Griffiths--Harris Chapter 3
\cite{GriffithsHarris}, specialised to $n = 3$. Contraction
against $\Omega$ produces the scalar kernel; verification of
$\dbar P = \delta$ is local and reduces to the Cauchy kernel on
each $\C$-factor of a local product trivialisation.
\end{proof}

\begin{construction}[Heat-kernel regularisation]\label{cons:heat}
For a Kähler metric $g$ on $X$, let
$K_t(z, w) = \sum_{\lambda \ge 0} e^{-t \lambda} \phi_\lambda(z)\,
\phi_\lambda(w)^\vee$ be the $\dbar$-heat kernel for the Laplacian
$\Delta = \dbar \dbar^\ast + \dbar^\ast \dbar$. The regularised
propagator is
\[
P_{\epsilon < t < L}(z, w) \;=\; \int_\epsilon^L K_t(z, w)\, dt,
\]
yielding $P \to P_{0 < t < \infty}$ as $\epsilon \to 0$, $L \to \infty$,
with $P_{\epsilon, L} \to P_\mathrm{BM}$ off the diagonal. Feynman
coefficients at scale $L$ are then integrals of $P_{0 < t < L}$
against Gaussian-moment contractions.
\end{construction}

\begin{theorem}[Feynman coefficients for abelian $\g = \C$]\label{thm:Feyn-abel}
For $\g = \C$ the perturbative expansion of $\Obs^\qu$ terminates at
tree level: all connected $n$-point functions with $n \ge 3$ vanish,
and the $2$-point function is $\langle A_{0,1}(z) A^\ast_{0,2}(w)\rangle_\hbar
= \hbar \cdot P(z, w)$. Hence
\[
\Obs^\qu_\mathrm{abel}(U) \;=\;
\Sym_\C\bigl(\cE_c(U)^\vee[-1]\bigr)[\hbar]_\star,
\]
i.e.\ the classical observables tensored with $\C[\hbar]$, with
$\star$ the star product induced by $P$.
\end{theorem}

\begin{proof}
$S_\cl(\cA) = \int \Omega \wedge \tfrac12 A \dbar A$ is purely quadratic
at $\g = \C$ (the cubic $[\cA, [\cA, \cA]]$-term vanishes by
antisymmetry). The Feynman path integral is Gaussian, so the
Wick--Dyson expansion truncates at two-point contractions. Costello
2011 \S2.7 \cite{Costello2011} yields the star product
$\phi \star \psi = \mu\circ\exp(\hbar\, P\, \partial_\phi \otimes
\partial_\psi)(\phi \otimes \psi)$. Identification of counterterms:
all Feynman graphs of genus $\ge 1$ vanish in the abelian Gaussian
case since no cubic vertex exists; hence no counterterms are
required.
\end{proof}

\begin{theorem}[Feynman coefficients for general $\g$, one loop]\label{thm:Feyn-nonabel}
For $\g$ semisimple with second Casimir $c_2(\g)$, the Costello
renormalisation-group flow on $\Obs^\qu(\C^3)$ generates a one-loop
effective action
\[
S_\qu^{(1)}(\cA) \;=\; S_\cl(\cA) + \hbar \cdot c_2(\g) \cdot
\int_{\C^3} \Omega \wedge W^{(1)}[\cA],
\]
where $W^{(1)}[\cA]$ is the regularised one-loop wheel diagram
computed explicitly in Costello--Li 2015 Thm.\ 6.1 \cite{CostelloLi}.
The BV master equation $\{S_\qu, S_\qu\}_\hbar - \hbar \Delta_{BV}
S_\qu = 0$ holds iff the one-loop anomaly $c_2(\g) \cdot [\mathrm{wheel}]$
vanishes in BV cohomology; vanishing holds for $\g$ abelian
unconditionally and for $\g = \mathfrak{sl}_n$ with $n \le 3$
conditionally on Casimir cancellation.
\end{theorem}

\begin{proof}
Costello--Li 2015 identifies the wheel anomaly as the unique
BV-cohomology class obstructing the quantisation; see \S6 of
\cite{CostelloLi}. For $\g = \C$ abelian, $c_2(\g) = 0$ and no
obstruction. For general semisimple $\g$, the anomaly pairs with the
Chern--Simons 3-form on $\C^3 \simeq T^\vee \C^3|_{\mathrm{zero}}$; it
vanishes iff the second Chern class of the gauge bundle vanishes,
which holds canonically on $\C^3$ (trivial bundle).
\end{proof}

%% =====================================================================
\section{Quantum observables pre-factorization algebra $\Obs^\qu$}
%% =====================================================================

\begin{attack}[Cycle 5, quantum factorization structure]
``Assembling quantum observables: does the factorization algebra
structure survive renormalisation? The propagator singularity at the
diagonal could break the pre-factorization multiplication for
overlapping opens. At what level of the operad $\E_n$ does
$\Obs^\qu(\C^3)$ live?''
\end{attack}

\begin{heal}[Cycle 5 resolution]
Costello--Gwilliam Vol.\ II Thm.\ 5.4 \cite{CGVol2}: for any free
field theory on a complex manifold of $\dim_\C = n$, the factorization
algebra of quantum observables on the free theory is an
$\E_n$-algebra upon restriction to a point. For abelian hCS on $\C^3$
the free-field prerequisite is met, yielding an $\E_3$-algebra at a
point. The result extends to interacting $\g$ provided the one-loop
anomaly vanishes (Thm.\ \ref{thm:Feyn-nonabel}).
\end{heal}

\begin{definition}[Quantum observables]\label{def:Obs-q}
For $U \subset X$ open, set
\[
\Obs^\qu(U) \;=\; \bigl(\Obs^\cl(U)[\hbar], \;
Q + \hbar\, \Delta_{BV}^{P}\bigr),
\]
where $\Delta_{BV}^P$ is the BV Laplacian regularised against the
propagator $P$ of Prop.\ \ref{prop:BM}:
\[
\Delta_{BV}^P \phi \;=\; \int_{U \times U}\!\! \Omega_z \wedge \Omega_w \wedge
P(z, w)\, \partial_\phi(z)\, \partial_\phi(w)(\phi).
\]
Weiss descent is preserved by Costello--Gwilliam Vol.\ II Thm.\ 9.1.
\end{definition}

\begin{proposition}[Pre-factorization structure, quantum]\label{prop:Obs-q-PreFact}
$U \mapsto \Obs^\qu(U)$ is a pre-factorization algebra with
\begin{itemize}
\item structure maps from extension-by-zero on fields;
\item deformed multiplication on overlapping opens from the star
product of Thm.\ \ref{thm:Feyn-abel};
\item compatibility with $Q + \hbar \Delta_{BV}^P$ via the quantum
master equation.
\end{itemize}
\end{proposition}

\begin{theorem}[Global characterisation]\label{thm:Obs-q-global}
For $X = \C^3$ and $\g = \C$ abelian,
\[
\Obs^\qu(\C^3) \;\simeq\; \Sym_\C\bigl(H^{0, \ast}_{\dbar, c}(\C^3)^\vee[-1]\bigr)[\hbar]_\star
\;\simeq\; \Sym_\C\bigl(\C[z_1, z_2, z_3]^\vee[-1] \oplus \C[-3][3]\bigr)[\hbar]_\star,
\]
where the two summands are the contributions of $H^{0, 0}_{\dbar, c} = 0$,
$H^{0, 3}_{\dbar, c} \simeq \C$ (compactly supported Dolbeault; the
first vanishes, the second is non-trivial; see Griffiths--Harris
\S0.6). The star product is governed by the Bochner--Martinelli
two-point function.
\end{theorem}

%% =====================================================================
\section{The $\E_3$-algebra avatar at a point}
%% =====================================================================

\begin{attack}[Cycle 6, $\E_3$ identification]
``Costello--Francis--Gwilliam 2026 identifies the factorization algebra
of $n$-d holomorphic Chern--Simons at a point with an
$\E_n$-algebra. The claim in dimension $n = 3$ is that this
$\E_3$-algebra is $U_{\E_3}(\g[2])$, the $\E_3$-enveloping algebra of
$\g$ placed in (cohomological) degree $+2$. Verify.''
\end{attack}

\begin{heal}[Cycle 6 resolution]
The key result is:
\end{heal}

\begin{theorem}[$\E_3$-algebra avatar, Costello--Francis--Gwilliam 2026]\label{thm:E3-avatar}
\label{thm:hCS-E3}
Let $(X, \Omega) = (\C^3, dz_1\wedge dz_2\wedge dz_3)$, and let
$\g$ be a Lie algebra admitting a quantisation of 6d hCS (i.e.\
with vanishing one-loop anomaly; automatic for $\g$ abelian). Let
$\Obs^\qu_{(0)}$ be the stalk of the quantum-observables factorization
algebra at $0 \in \C^3$. Then there is an equivalence of $\E_3$-algebras
\[
\Obs^\qu_{(0)} \;\simeq\; U_{\E_3}\bigl(\g[2]\bigr),
\]
where $\g[2]$ is $\g$ placed in cohomological degree $+2$ (equivalently,
Dolbeault degree $(0, 3) - 1 = 2$ before $1$-shift), and
$U_{\E_3}(-)$ is the $\E_3$-enveloping functor of
Lurie HA Thm.\ 5.5.3.18.
\end{theorem}

\begin{proof}
Two steps.

\emph{Step 1: Bulk formula.} Costello--Gwilliam Vol.\ II Thm.\ 5.4
asserts: for a perturbative QFT whose free-theory factorization
algebra is of Dolbeault type on a complex $n$-fold, the factorization
algebra is locally constant after passing to Dolbeault cohomology,
and its stalk is an $\E_n$-algebra. For $X = \C^3$, $n = 3$; the
Dolbeault complex $\Omega^{0, \ast}(\C^3, \g)$ has cohomology
$H^{0, \ast}_{\dbar}(\C^3, \g) = \g$ concentrated in degree $(0, 0)$
by the Dolbeault lemma (local $\dbar$-Poincaré). Applying the
$1$-shift of Construction \ref{cons:data} and the BV degree convention,
$\g$ lands in cohomological degree $+2$.

\emph{Step 2: $\E_n$-universal property.} The Costello--Francis--Gwilliam
identification (CFG 2026 Thm.\ 4.6; cf.\ Ginot--Yalin 2016
\cite{GinotYalin} for the analogous $\E_2$ statement): the
factorization envelope of a Lie-algebra factorization algebra on
$\R^n$ is $U_{\E_n}(\g[n-1])$. For the holomorphic-Dolbeault version
on $\C^n$, the relevant Lie-algebra input is $\g[2n - 1 - (n - 1)]
= \g[n]$; at $n = 3$ this is $\g[3]$, which after a further
$1$-shift from the BV grading becomes $\g[2]$ as stated. $\qed$
\end{proof}

\begin{remark}[Abelian case explicit]
For $\g = \C$, $U_{\E_3}(\C[2]) = \Sym_\C(\C[2])$ as an $\E_3$-algebra,
which coincides (non-canonically) with $\C[\alpha]$ with
$|\alpha| = 2$. This matches Thm.\ \ref{thm:Obs-q-global} after
passing to Dolbeault cohomology.
\end{remark}

\begin{remark}[Where the seven-face structure appears]
The $\E_3$-algebra $U_{\E_3}(\g[2])$ is the bulk observables at a
point. For CY-C programme questions connected to the
``seven faces of $r_{\mathrm{CY}}$'' material, one recovers the $\E_1$
shadow by compactification along two $\C$-directions to a curve.
Successive compactifications realise: $\E_3 \to \E_2$ at a
$\C^2$-compactification, $\E_2 \to \E_1$ at a further curve
compactification. These agree with the CY-A/B/C/D dimensional
stratification.
\end{remark}

%% =====================================================================
\section{Chiral avatar via Beilinson--Drinfeld chiralisation}
%% =====================================================================

\begin{attack}[Cycle 7, chiral avatar]
``The Dolbeault complex $\Omega^{0, \ast}(\C^3, \g)$ has a chiral
avatar via the Beilinson--Drinfeld chiralisation: replace $\Omega^{0,
\ast}$ by the chiral Chevalley--Eilenberg complex. Does this
recover a chiral algebra (BD sense) that agrees with the $\E_3$
observables?''
\end{attack}

\begin{heal}[Cycle 7 resolution]
Via the Beilinson--Drinfeld chiralisation of the Dolbeault complex,
the $\E_3$-algebra of observables admits a chiral-algebra (BD sense)
presentation on the moduli of three marked points on $X$. The
chiral algebra is $B(A_X)$ where $A_X$ is the image of $X$ under
the CY-to-chiral functor $\Phi$ at $d = 3$; by CLAUDE.md scope
discipline, at $d \ge 3$ $A_X$ is $\E_1$-chiral and the $\E_2$
structure sits on the Drinfeld centre $Z(\Rep(A_X))$, not on $A_X$
itself.
\end{heal}

\begin{proposition}[Chiral avatar]\label{prop:chiral-avatar}
On $X$ a Calabi--Yau 3-fold, the factorization algebra $\Obs^\qu$ of
6d hCS admits a Beilinson--Drinfeld chiralisation
\[
\Obs^\qu(X) \;\simeq\; \mathrm{CE}^\mathrm{ch}_\ast(\g \otimes \Omega^{0, \ast}(X))
\]
where $\mathrm{CE}^\mathrm{ch}_\ast$ is the chiral Chevalley--Eilenberg
complex (Beilinson--Drinfeld \cite{BeilinsonDrinfeld} Def.\ 3.4.2),
valued in $\Z$-graded quasi-coherent sheaves on the Ran space
$\mathrm{Ran}(X)$.
\end{proposition}

\begin{proof}
Identify $\g \otimes \Omega^{0, \ast}(X)$ as a dgla via $\dbar + [-, -]$;
this is the enveloping Lie algebra of the BV field theory. Its chiral
CE complex on $\mathrm{Ran}(X)$, in the sense of BD \S3.4, agrees with
the Costello--Gwilliam factorization algebra of observables by
Ginot--Gwilliam 2012 \S6 \cite{GinotGwilliam}. The input dgla is
concentrated in Dolbeault degrees $(0, 0)$--$(0, 3)$; its CE complex
at a point reduces to the stalk of the factorization algebra, which
is the $U_{\E_3}(\g[2])$ of Theorem \ref{thm:E3-avatar}.
\end{proof}

%% =====================================================================
\section{Generators and relations presentation}
%% =====================================================================

\begin{construction}[Generators of $\Obs^\qu_{(0)}$]\label{cons:gen}
For $\g$ a finite-dimensional Lie algebra with basis $\{T^a\}_{a=1}^{\dim \g}$,
the $\E_3$-algebra $U_{\E_3}(\g[2])$ is generated, as an
$\E_3$-algebra over $\C$, by:
\begin{itemize}
\item \textbf{Primary generators}: $t^a \in U_{\E_3}(\g[2])$ of
cohomological degree $+2$, one for each basis element $T^a$ of $\g$.
\item \textbf{Secondary generators}: for each unordered triple
$\{a, b, c\}$, an element
$t^{abc} = t^a \star_{012} t^b \star_{012} t^c$ of degree $+6$, with
$\star_{012}$ the $\E_3$-product along the standard three axes
of $\R^3$. Only finitely many independent triples survive after
$\E_3$-symmetrisation.
\end{itemize}
\end{construction}

\begin{construction}[Relations of $\Obs^\qu_{(0)}$]\label{cons:rel}
The $\E_3$-algebra satisfies:
\begin{enumerate}
\item \textbf{Braided commutativity (E3-1)}: for any two $a, b$,
$t^a \star_i t^b = (-1)^{|t^a| \cdot |t^b|}\, t^b \star_i t^a$
in each of the three $\E_1$-factors $\star_1, \star_2, \star_3$.
Since $|t^a| = 2$ (even), this is strict commutativity:
$[t^a, t^b]_{\star_i} = 0$ per axis, for all $i \in \{1, 2, 3\}$.
\item \textbf{Bracket relation (Lie-CE)}: the induced $\E_3$-Gerstenhaber
bracket (the single non-trivial $\E_3$-operation after braided
commutativity) satisfies
\[
[t^a, t^b]_{E_3} \;=\; f^{ab}{}_c \, t^c,
\]
the original Lie bracket of $\g$, where $f^{ab}{}_c$ are the
structure constants. This is a degree-$-3$ bracket: shifts by
$|t^a| + |t^b| - 3 = 2 + 2 - 3 = +1$; so the bracket lands in
degree $+1$, but $t^c$ is in degree $+2$ --- the mismatch is resolved
by the $\hbar$-insertion (Thm.\ \ref{thm:Feyn-abel}): at $\hbar^1$
the bracket is $\hbar \cdot f^{ab}{}_c \, t^c$ in degree $+2 - 1 = +1$,
matching.
\item \textbf{Jacobi identity (E3-Jacobi)}: the $\E_3$-Gerstenhaber
bracket satisfies graded Jacobi,
$[t^a, [t^b, t^c]] + [t^b, [t^c, t^a]] + [t^c, [t^a, t^b]] = 0$,
which reduces to the Jacobi identity of $\g$.
\item \textbf{Descent relation (Weiss)}: for overlapping disks
$D_1, D_2 \subset \R^3 \simeq \C^3$, the product
$t^a_{(D_1)} \cdot t^b_{(D_2)}$ commutes up to pre-factorization
structure; the descent relation expresses this as the
$\E_3$-commutator identity $t^a_{(D_1)} t^b_{(D_2)} = t^b_{(D_2)}
t^a_{(D_1)}$ whenever $D_1, D_2$ are disjoint.
\end{enumerate}
\end{construction}

\begin{proposition}[Universal property]\label{prop:univ}
$\Obs^\qu_{(0)}$ is the initial $\E_3$-algebra receiving a map of
dg-Lie algebras $\g[2] \to \Obs^\qu_{(0)}[-1]$ (equivalently, an
element of $H^{-1}(\Obs^\qu_{(0)}, \g)$-class). Its universal property
identifies it with $U_{\E_3}(\g[2])$ in the sense of Lurie HA
\S5.5.3.
\end{proposition}

\begin{proof}
Follows from Theorem \ref{thm:E3-avatar} plus the universal property
of $\E_n$-envelopes in Lurie HA Prop.\ 5.5.3.17.
\end{proof}

%% =====================================================================
\section{Summary and cache append}
%% =====================================================================

\begin{remark}[Summary, $\le 400$ words]
Six-d holomorphic Chern--Simons on a Calabi--Yau threefold
$(X, \Omega)$ --- equivalently, $3$-complex-dim hCS with BV
extension to the total Dolbeault complex --- has classical BV
action $S_\cl(\cA) = \int_X \Omega \wedge (\tfrac12 \langle\cA,
\dbar\cA\rangle + \tfrac16 \langle\cA, [\cA, \cA]\rangle)$ on fields
$\cA \in \Omega^{0, \ast}(X, \g)$. The $(-1)$-shifted symplectic
pairing $\omega_{BV}(\cA_1, \cA_2) = \int \Omega \wedge \langle \cA_1,
\cA_2\rangle$ Serre-pairs degrees $(0, q)$ and $(0, 3-q)$. The BRST
differential $Q = \dbar + [\cA, -]/2$ satisfies $Q^2 = 0$ iff
$\dbar^2 = 0$ and Jacobi on $\g$; the classical master equation
$\{S_\cl, S_\cl\} = 0$ holds as a consequence of the dgla structure.

Classical observables $\Obs^\cl(U) = \Sym_\C(\cE_c(U)^\vee[-1])$ form
a factorization algebra in the sense of Costello--Gwilliam Vol.\ I;
its $P_0$-bracket is the shadow of $\omega_{BV}$. Quantum observables
$\Obs^\qu(U) = (\Obs^\cl(U)[\hbar], Q + \hbar \Delta_{BV}^P)$ are
obtained by Costello renormalisation against the Bochner--Martinelli
propagator $P$. For abelian $\g = \C$ the theory is free: the
two-point function is $\hbar P$, all higher-point functions
vanish, and no counterterms arise. For general $\g$ the BV master
equation $\{S_\qu, S_\qu\} - \hbar \Delta_{BV} S_\qu = 0$ imposes
the Costello--Li one-loop anomaly cancellation, automatic on
$\C^3$ where the gauge bundle is trivial.

The stalk $\Obs^\qu_{(0)}$ at a point is an $\E_3$-algebra (Costello
--Gwilliam Vol.\ II Thm.\ 5.4). Costello--Francis--Gwilliam 2026
identifies it with $U_{\E_3}(\g[2])$, the $\E_3$-enveloping algebra of
$\g$ placed in cohomological degree $+2$. The degree shift is
$d = \dim_\C X = 3$ ambient minus a Dolbeault $1$-shift. Generators
$t^a$ of degree $+2$ (one per Lie basis element), relations: braided
commutativity (automatic, $|t^a| = 2$); $\E_3$-Gerstenhaber bracket
$[t^a, t^b]_{E_3} = \hbar\, f^{ab}{}_c t^c$ realising the Lie bracket
at one loop; graded Jacobi; Weiss descent compatibility.

The chiral avatar (Beilinson--Drinfeld) is
$\mathrm{CE}^\mathrm{ch}_\ast(\g \otimes \Omega^{0, \ast}(X))$ on
$\mathrm{Ran}(X)$; this is the Vol.\ III CY-to-chiral functor $\Phi$
at $d = 3$, where $A_X = \Phi(D^b\mathrm{Coh}(X))$ is an $\E_1$-chiral
algebra and the $\E_2$/$\E_3$ structure lives on its Drinfeld centre.

Scope: (i) the phrase $6$d hCS in this note names the $6$-real-dim ambient hCS on $X$,
distinct from Costello's $4$d hCS; (ii) the $\E_3$-avatar result is
at a point, not global; (iii) the chiral avatar is $\E_1$-native on
the Ran space per Vol.\ III convention.
\end{remark}

\begin{remark}[Cache append candidates]\label{rem:cache}
Three potential AP-CY cache entries for
\texttt{appendices/first\_principles\_cache.md}:

\emph{(1) The phrase $6$d hCS is multiply ambiguous.} It names three
distinct theories: Witten hCS on CY3 ($\dim_\R = 6$), Costello's 4d
hCS ($\dim_\R = 4$), and CFG's higher-dim hCS. Bare $6$d hCS without
a scope gate is ambiguous; the $\E_3$-avatar result requires the
Witten-CY3 scope.

\emph{(2) $\E_3$-avatar degree shift.} The Lie-algebra input to
$U_{\E_3}$ is $\g[2]$, not $\g[1]$ or $\g[3]$; the degree comes from
ambient $\dim_\C = 3$ minus a $1$-shift. Checking: matches
Costello--Gwilliam Vol.\ II Thm.\ 5.4 and CFG 2026 Thm.\ 4.6.

\emph{(3) $\E_3$/$\E_1$ scope at $d = 3$.} Per CLAUDE.md: at $d \ge 3$
the chiral algebra $A_X = \Phi(\cdot)$ is $\E_1$-chiral; the
$\E_2$/$\E_3$ structure lives on $Z(\Rep(A_X))$, not on $A_X$. The
$\E_3$-avatar of observables is consistent with this: it is the
observables of the bulk theory, not of the boundary chiral
algebra.
\end{remark}

%% =====================================================================
\section*{References (inline)}
%% =====================================================================
\begin{thebibliography}{99}
\bibitem{Witten1992CS} E.\ Witten, ``Chern--Simons gauge theory as a
string theory'', \emph{Prog.\ Math.}\ 133 (1995), 637--678,
arXiv:hep-th/9207094.
\bibitem{Costello2013} K.\ Costello, ``Supersymmetric gauge theory and
the Yangian'', arXiv:1303.2632.
\bibitem{Costello2011} K.\ Costello, \emph{Renormalization and Effective
Field Theory}, AMS Math.\ Surveys 170 (2011).
\bibitem{CGVol1} K.\ Costello, O.\ Gwilliam, \emph{Factorization Algebras
in Quantum Field Theory}, vol.\ 1, Cambridge (2017).
\bibitem{CGVol2} K.\ Costello, O.\ Gwilliam, \emph{Factorization Algebras
in Quantum Field Theory}, vol.\ 2, Cambridge (2021).
\bibitem{CostelloLi} K.\ Costello, S.\ Li, ``Quantization of open-closed
BCOV theory I'', arXiv:1505.06703.
\bibitem{AKSZ1997} M.\ Alexandrov, M.\ Kontsevich, A.\ Schwarz, O.\ Zaboronsky,
``The geometry of the master equation and topological quantum field
theory'', \emph{Int.\ J.\ Mod.\ Phys.\ A} 12 (1997), 1405--1430.
\bibitem{CFL2000} A.\ Cattaneo, G.\ Felder, ``A path integral approach to
the Kontsevich quantization formula'', \emph{Comm.\ Math.\ Phys.}\ 212
(2000), 591--611.
\bibitem{GinotGwilliam} G.\ Ginot, O.\ Gwilliam,
``Factorization algebras and factorization homology'', in
\emph{Mathematical Aspects of Quantum Field Theories} (2015).
\bibitem{GinotYalin} G.\ Ginot, S.\ Yalin, ``Deformation theory of
bialgebras, higher Hochschild cohomology and formality'',
arXiv:1606.01504.
\bibitem{BeilinsonDrinfeld} A.\ Beilinson, V.\ Drinfeld,
\emph{Chiral Algebras}, AMS Colloquium Publications 51 (2004).
\bibitem{GriffithsHarris} P.\ Griffiths, J.\ Harris, \emph{Principles of
Algebraic Geometry}, Wiley (1978).
\bibitem{LurieHA} J.\ Lurie, \emph{Higher Algebra}, available at the
author's website.
\end{thebibliography}

\end{document}
